\BourbakiExerciseGroup

\begin{enumerate}
\def\labelenumi{\arabic{enumi})}
\tightlist
\item
  设 \(\mathbf{K}\) 是特征为 \(p > 0\) 的域，\(\mathbf{E}\) 是 \(\mathbf{K}\) 的一个扩张，\(\mathbf{B}\) 是一个 \(p\) -基（\(\mathbf{E}\) 在 \(\mathbf{K}(\mathbf{E}^p)\) 上的）（V，第 98 页）。
\end{enumerate}

\begin{enumerate}
\def\labelenumi{\alph{enumi})}
\item
  假设 \(\mathbf{E} \subset \mathbf{K}^{p^{-1}}\) ；设 \(\mathbf{K}_0\) 是 \(\mathbf{K}\) 的一个子域，使得 \(\mathbf{E}\) 在 \(\mathbf{K}_0\) 上可分。证明集合 \(\mathbf{B}^p\) 是一个 \(p\) -无关子集（在 \(\mathbf{K}_0(\mathbf{K}^p)\) 上）（注意，若 \((a_{\lambda})\) 是向量 \(K_{0}\) -空间 K 的一组基，则 \((a_{\lambda}^{p})\) 是 \(K_{0}(K^{p})\) 在 \(K_{0}\) 上的一组基。设 C 是 K 的一个子集，与 \(B^{p}\) 不相交，且使 \(B^{p} \cup C\) 是 K 在 \(K_{0}(K^{p})\) 上的一个 p-基；证明 \(B \cup C\) 是 E 在 \(K_{0}(E^{p})\) 上的一个 p-基。
\item
  假设 \(E \subset K^{p^{-n}}\) ，且 E 在 K 的一个子域 \(K_{0}\) 上可分。证明：若 K 在 \(K_{0}\) 上的不完全度有限（V，第 170 页，习题 1），则它等于 E 在 \(K_{0}\) 上的不完全度（利用 a））。
\end{enumerate}

\begin{enumerate}
\def\labelenumi{\arabic{enumi})}
\setcounter{enumi}{1}
\tightlist
\item
  设 \(\mathbf{K}\) 是特征为 \(p > 0\) 的域，\(\mathbf{E}\) 是 \(\mathbf{K}\) 的一个扩张，\(\mathbf{F}\) 是 \(\mathbf{E}\) 的一个可分扩张。证明下列三个性质中的任意两个蕴含第三个：
\end{enumerate}

\begin{enumerate}
\def\labelenumi{\alph{enumi})}
\item
  B 是一个 \(p\) -基（E 在 \(\mathbf{K}(\mathbf{E}^p)\) 上的）；
\item
  C 是一个 \(p\) -基（\(\mathbf{F}\) 在 \(\operatorname{E}(\mathbf{F}^p)\) 上的）；
\item
  \(\mathbf{B} \subset \mathbf{E}\) ，\(\mathbf{B} \cup \mathbf{C}\) 是一个 \(p\) -基（\(\mathbf{F}\) 在 \(\mathbf{K}(\mathbf{F}^p)\) 上的），且 \(\mathbf{B} \cap \mathbf{C} = \emptyset\) 。（利用如下事实：若 \((c_{\mu})\) 是向量 E-空间 \(\mathbf{F}\) 的一组基，则 \((c_{\mu}^p)\) 是 \(\mathbf{K}(\mathbf{E}^p)[\mathbf{F}^p]\) 在 \(\mathbf{K}(\mathbf{E}^p)\) 上以及 \(\mathbf{E}[\mathbf{F}^p]\) 在 \(\mathbf{E}\) 上的一组基。）
\end{enumerate}

\begin{enumerate}
\def\labelenumi{\arabic{enumi})}
\setcounter{enumi}{2}
\item
  设 K 是特征为 p \textgreater{} 0 的域，E 是 K 的一个扩张，F 是 E 的一个有限生成扩张。证明 F 在 K 上的不完全度至少等于 E 在 K 上的不完全度（归结为以下两种情形： \(1^{\circ}\) F 在 E 上可分； \(2^{\circ}\) F = E(x)，其中 \(x^{p} \in E\) （参见 V，第 170 页，习题 2））。
\item
  设 \(\mathbf{F}\) 是域 \(\mathbf{K}\) —— 特征为 \(p > 0\) —— 的一个可分扩张。证明：若 \(\mathbf{E}\) 是 \(\mathbf{F}\) 的一个在 \(\mathbf{K}\) 上相对完全的子扩张（V，第 170 页，习题 1），则 \(\mathbf{F}\) 在 \(\mathbf{E}\) 上可分。
\item
  设 \(\mathbf{K}\) 是特征为 \(p > 0\) 的域。证明：若 \(\mathbf{E}\) 是 \(\mathbf{K}\) 的代数扩张或 \(\mathbf{K}\) 的相对完全扩张，则 \(\mathrm{E}^{p^{-\infty}} = \mathrm{E}(\mathrm{K}^{p^{-\infty}})\) 。
\item
  设 \(\mathbf{K}\) 是特征为 \(p > 0\) 的域，\(\mathbf{E}\) 是 \(\mathbf{K}\) 的一个可分扩张，\(\mathbf{B}\) 是一个 \(p\) -基（\(\mathbf{E}\) 在 \(\mathbf{K}(\mathbf{E}^p)\) 上的）。
\end{enumerate}

\begin{enumerate}
\def\labelenumi{\alph{enumi})}
\item
  证明 B 在 K 上是代数自由的（考虑 B 的元素之间次数最低的一个代数关系，并将其中出现的变元的次数写成 \(kp + h\) 的形式，其中 \(0 \leqslant h \leqslant p - 1\) ）。由此推出：若 tr. \(\deg_{K}E < +\infty\) ，则 E 在 K 上的不完全度至多等于 tr. \(\deg_{K}E\) 。
\item
  证明 E 在 K(B) 上是可分的且相对完美的。
\end{enumerate}

\begin{enumerate}
\def\labelenumi{\arabic{enumi})}
\setcounter{enumi}{6}
\item
  证明特征为 p \textgreater{} 0 的域 K 的一个相对完全超越扩张 E 在 K 上不是有限生成的。（在相反的情形，证明对 E 在 K 上的每个超越基 B，E 都在 K(B) 上是代数的且可分的。）
\item
  若 E 是特征为 p \textgreater{} 0 的域 K 的一个可分代数扩张，则 E 的最大完全子扩张 \(E^{p^{\infty}}\) （诸域 \(E^{p^{n}}\) （ \(n \geqslant 0\) ）的交）在 K 的最大完全子域 \(K^{p^{\infty}}\) 上是代数的。（若 \(x \in E^{p^{\infty}}\) ，证明对每个整数 n \textgreater{} 0，我们有 \(x^{p^{-n}} \in \mathrm{K}(x)\) ，并由此推出 x 在 \(K^{p^{n}}\) 上的极小多项式与它在 K 上的极小多项式相同。）
\item
  设 \(\mathbf{E}\) 是一个域，\(\mathbf{G}\) 是 \(\mathbf{E}\) 的一个自同构群，而 \(\mathbf{K}\) 是 \(\mathbf{E}\) 的由在 \(\mathbf{G}\) 下不变的元素组成的子域。
\end{enumerate}

\begin{enumerate}
\def\labelenumi{\alph{enumi})}
\item
  证明：对 E 的每个在 G 的元素下稳定的子域 L，L 与 K 在 L ∩ K 上线性不相交。（考虑一个线性关系 \(\sum_{j}\lambda_{j}x_{j}=0\) ，其中的元素 \(x_{j} \in \mathbf{K}\) 在 \(\mathbf{L} \cap \mathbf{K}\) 上线性无关，而 \(\lambda_{j} \in \mathbf{L}\) 不全为零，且 \(\lambda_{j} \neq 0\) 的个数尽可能地少。）
\item
  若 \(L \cap K\) 是 K 的某个自同构群的不变域，证明 L 是 \(\mathrm{L}(\mathbf{K})\) 的某个自同构群的不变域。
\end{enumerate}

\begin{enumerate}
\def\labelenumi{\arabic{enumi})}
\setcounter{enumi}{9}
\item
  设 E 是域 K 的一个扩张，G 是拓扑群 \(\operatorname{Aut}_{\mathbf{K}}(\mathbf{E})\) 的一个紧子群，F 是 G 的不变域。证明 E 是 F 的一个伽罗瓦扩张，且 G 典范地同构于拓扑群 \(\operatorname{Aut}_{\mathbf{F}}(\mathbf{E})\) 。（注意 G 以连续方式作用于离散空间 E，以确立元素 \(x \in E\) 在 G 的作用下具有有限轨道。）
\item
  设 \(\mathbf{K}\) 是特征为 \(p > 0\) 的域。对每个交换 K-代数 \(\mathbf{A}\) ，取 \(p\) 次幂使环 \(\mathbf{A}\) 成为一个 A-代数，从而也是一个 K-代数，记作 \(\mathbf{A}^{p^{-1}}\) 。设 \(\mathbf{A}\) 是一个交换 K-代数。证明下列条件是等价的：
\end{enumerate}

\begin{enumerate}
\def\labelenumi{(\roman{enumi})}
\item
  \(\mathbf{A}\) 是一个可分的 K-代数。
\item
  唯一的 K-同态 \(\Phi: \mathbf{A} \otimes_{\mathbf{K}} \mathbf{K}^{p^{-1}} \to \mathbf{A}^{p^{-1}}\) 使得 \(\Phi(a \otimes \lambda) = a^p \cdot \lambda (a \in \mathbf{A}, \lambda \in \mathbf{K})\) 是单射的。
\item
  对每个族 \((b_{i})_{i\in I}\) （在 \(\mathbf{K}\) 中、于 \(\mathbf{K}^p\) 上线性无关）以及每个族 \((a_{i})_{i\in I}\) （在 \(\mathbf{A}\) 中），等式 \(\sum_{i}a_i^p b_i = 0\) 蕴含 \(a_{i} = 0\) 对所有 \(i\) 成立。
\end{enumerate}

（可以验证（ii） \(\Leftrightarrow\) （iii），并且（ii）是V，第123页，定理2，e）的推论。）

\begin{enumerate}
\def\labelenumi{\arabic{enumi})}
\setcounter{enumi}{11}
\tightlist
\item
  设 K 是一个域，\((\mathbf{X}_{i})_{i\in I}\) 是一族未定元。证明形式幂级数代数 \(\mathbf{K}[[(\mathbf{X}_{i})_{i\in I}]]\) 是一个可分 K-代数。（应用习题 11；也可以验证，对 K 的每个有限扩张 \(K'\) ，\(\mathbf{K}[[(\mathbf{X}_{i})_{i\in I}]]\otimes_{\mathbf{K}}K'\) 同构于 \(\mathbf{K}'[[(\mathbf{X}_{i})_{i\in I}]]\) ，然后应用 V，第 123 页，定理 2，d)。）由此推出 \(\mathbf{K}[[(\mathbf{X}_{i})_{i\in I}]]\) 的分式域是 K 的一个可分扩张（V，第 120 页，命题 4）。
\end{enumerate}

